%%
%% ACM TOSEM — Adapted for: A Taxonomy of Code Repository Representation
%% Approaches for LLM-Based Software Engineering: A Systematic Literature Review
%%
\documentclass[acmsmall,screen,review]{acmart}

\AtBeginDocument{%
  \providecommand\BibTeX{{%
    \normalfont B\kern-0.5em{\scshape i\kern-0.25em b}\kern-0.8em\TeX}}}

%% Rights management
\setcopyright{acmcopyright}
\copyrightyear{2026}
\acmYear{2026}
\acmDOI{XXXXXXX.XXXXXXX}

\acmJournal{TOSEM}
\acmVolume{35}
\acmNumber{1}
\acmArticle{1}
\acmMonth{4}

\usepackage{booktabs}
\usepackage{multirow}
\usepackage{array}
\usepackage{rotating}
\usepackage{xcolor}
\usepackage{amsmath}
\usepackage{bm}
\usepackage[most]{tcolorbox}

\begin{document}

\title{A Taxonomy of Code Repository Representation Approaches for
  LLM-Based Software Engineering: A Systematic Literature Review}

\author{Riadh Hasnaoui}
\email{riadh.hasnaoui@ensi-uma.tn}
\affiliation{%
  \institution{National School of Computer Science (ENSI)}
  \city{Manouba}
  \country{Tunisia}
}

%%
%% Abstract
%%
\begin{abstract}
Large Language Models (LLMs) have demonstrated exceptional capabilities
in software engineering (SE) tasks, yet effectively providing entire
code repositories as context remains a critical open challenge. This
Systematic Literature Review (SLR) proposes a comprehensive, empirically
grounded taxonomy of repository representation approaches for LLM-based
SE. Following a rigorous \textbf{five-pillar methodological framework}
(Kitchenham \& Charters, Wohlin, PRISMA 2020, Cruzes \& Dyb\aa{}, and
ACM SIGSOFT Empirical Standards), we systematically analysed a corpus
of \textbf{159 studies} (134 primary representation studies and 25
benchmark evaluation studies) published between January 2020 and
March 2026.

Our thematic synthesis identifies \textbf{seven distinct representation
paradigms (PAR-1 to PAR-7)}: Dense Embedding, RAG Cross-File,
Graph-Based, Execution \& Dynamic, Memory-Augmented, Compression \&
Summarisation, and Agentic \& Iterative Context Exploration. Each
paradigm is characterised across \textbf{ten structural dimensions
(D1–D10)}, forming the taxonomy's analytical backbone. Our
cross-paradigm analysis reveals (1) a fundamental trade-off between
token-budget consumption and semantic precision; (2) a systematic
\emph{freshness bias} in four paradigms that rely on stale or
semi-stale repository snapshots, leading to performance overestimation
on static benchmarks; and (3) an unresolved \emph{training-dependency
gap} that limits the transferability of encoder-centric paradigms to new
programming languages. This SLR provides a foundational taxonomy and
research agenda to guide future research and the industrial deployment
of repository-level LLM agents.
\end{abstract}

\begin{CCSXML}
<ccs2012>
 <concept>
  <concept_id>10011007.10011074.10011092</concept_id>
  <concept_desc>Software and its engineering~Software development
   techniques</concept_desc>
  <concept_significance>500</concept_significance>
 </concept>
 <concept>
  <concept_id>10010147.10010178.10010179</concept_id>
  <concept_desc>Computing methodologies~Natural language
   processing</concept_desc>
  <concept_significance>500</concept_significance>
 </concept>
 <concept>
  <concept_id>10011007.10011074.10011099.10011102.10011103</concept_id>
  <concept_desc>Software and its engineering~Software testing and
   debugging</concept_desc>
  <concept_significance>300</concept_significance>
 </concept>
</ccs2012>
\end{CCSXML}

\ccsdesc[500]{Software and its engineering~Software development techniques}
\ccsdesc[500]{Computing methodologies~Natural language processing}
\ccsdesc[300]{Software and its engineering~Software testing and debugging}

\keywords{Large Language Models, Software Engineering, Systematic
  Literature Review, Repository Representation, Code Generation, RAG,
  Autonomous Agents, Taxonomy, Code Completion, Program Repair}

\maketitle

%%
%% ─── SECTION 1: INTRODUCTION ────────────────────────────────────────────────
%%
\section{Introduction}

The integration of Large Language Models (LLMs) into Software
Engineering (SE) has fundamentally shifted the paradigm from
function-level code assistance to repository-level
automation~\cite{Hou2024, Jiang2026, Zhang2026Survey}. Systems such as
SWE-agent~\cite{Yang2024SWEAgent}, AutoCodeRover~\cite{Zhang2024ACR},
and OpenHands~\cite{Wang2025OpenHands} demonstrate that LLMs can, in
principle, autonomously navigate entire codebases, localise faults, and
generate verified patches. Yet the core engineering challenge that
underlies all these systems remains largely unsolved: \emph{how does one
represent a complex, multi-file repository so that an LLM can understand
its cross-file dependencies and structural constraints without exceeding
the model's token limit?}

The question is not merely practical. The choice of representation
strategy determines which semantic features are visible to the model,
which are silently discarded, and which are implicitly approximated. As
the SWE-bench~\cite{Jimenez2024} benchmark has shown, even the most
capable proprietary models fail to resolve the majority of real-world
GitHub issues when repository context is naïvely truncated. Recent work
has confirmed that performance on static benchmarks can be
systematically overestimated when the repository snapshot used for
evaluation predates the model's training
cut-off~\cite{Zheng2025HumanEvo}, raising serious questions about the
validity of the current evaluation ecosystem.

Despite the rapid proliferation of repository representation techniques
-- from dense-vector retrieval~\cite{Feng2020CodeBERT, Guo2022UniX} to
graph-based structural
encoding~\cite{Guo2021GraphCodeBERT, Liu2025CodexGraph}, from
compression-based
sketching~\cite{Zan2025CodeS, Shi2025LongCodeZip} to agentic
exploration~\cite{Yang2024SWEAgent, Zhang2024ACR} -- no unified,
theoretically grounded taxonomy currently exists. Existing surveys
either focus on a single technique family (e.g., RAG for
code~\cite{Tao2026RAGSurvey}) or treat representation as a secondary
concern within broader LLM-for-SE surveys~\cite{Hou2024, Husein2025}.

\textbf{Contributions.} This paper makes the following contributions:
\begin{description}
  \item[C1 (Taxonomy)] A seven-paradigm taxonomy of repository
    representation approaches, characterised across ten structural
    dimensions (Section~\ref{sec:taxonomy}).
  \item[C2 (Corpus)] A systematically assembled corpus of 159 studies
    screened according to a five-pillar SLR protocol
    (Section~\ref{sec:methodology}).
  \item[C3 (Analysis)] A quantitative cross-paradigm analysis that
    exposes three fundamental trade-offs — token budget vs.\ semantic
    precision, freshness bias, and training dependency — and maps them
    to seven SE objectives (Sections~\ref{sec:taxonomy}
    and~\ref{sec:discussion}).
  \item[C4 (Agenda)] A structured research agenda for the field, derived
    from the identified gaps (Section~\ref{sec:agenda}).
\end{description}

\textbf{Research Questions.}
\begin{description}
  \item[RQ1] What are the distinct paradigms of repository representation
    for LLM-based SE, and how do they differ across structural
    dimensions?
  \item[RQ2] How are these paradigms mapped to automated SE objectives,
    and what are the empirically observed performance trade-offs?
  \item[RQ3] What are the most critical open challenges and directions
    for future research?
\end{description}

%%
%% ─── SECTION 2: METHODOLOGY ─────────────────────────────────────────────────
%%
\section{Research Methodology}
\label{sec:methodology}

This SLR was conducted according to a \textbf{five-pillar
methodological framework} that combines the rigour of established SLR
protocols with the analytical depth required to produce a taxonomy of
scientific validity.

\subsection{Pillar 1 — SLR Protocol: Kitchenham \& Charters}

We adopt the Kitchenham \& Charters~\cite{Kitchenham2007} guidelines as
our foundational protocol. The review scope was formalised through the
PICOC framework: \emph{Population} = code repositories processed as LLM
context; \emph{Intervention} = representation method; \emph{Comparison}
= alternative methods on the same SE task; \emph{Outcome} = accuracy,
token efficiency, TRL; \emph{Context} = automated SE.

\noindent\textbf{Inclusion Criteria (IC):}
\begin{itemize}
  \item IC1: The study proposes or empirically evaluates a method for
    extracting or structuring repository-level context for LLMs.
  \item IC2: The method targets at least one of the seven SE objectives
    (code completion, bug localisation, bug fixing, code search,
    program repair, code translation, test generation).
  \item IC3: The study is a peer-reviewed publication or an arXiv
    preprint with $\geq 10$ citations or $\geq 1$ replication at the
    time of screening.
\end{itemize}

\noindent\textbf{Exclusion Criteria (EC):}
\begin{itemize}
  \item EC1: Studies that operate exclusively at function or file level
    without cross-file or repository-level context.
  \item EC2: Studies focused solely on model architecture improvements
    without any repository representation contribution.
  \item EC3: Duplicate publications (we retained the most complete
    version).
\end{itemize}

\subsection{Pillar 2 — Snowballing: Wohlin (2014)}

Given the extreme publication velocity in LLM-for-code research, we
applied Wohlin's~\cite{Wohlin2014} snowballing protocol as a
completeness safeguard. Backward snowballing from the initial corpus
captured foundational works such as CodeBERT~\cite{Feng2020CodeBERT}
and GraphCodeBERT~\cite{Guo2021GraphCodeBERT}. Forward snowballing via
the Semantic Scholar API captured recent studies. Snowballing
contributed \textbf{27 studies} (17\% of the primary corpus).

\subsection{Pillar 3 — PRISMA 2020: Transparent Reporting}

This SLR adheres to all 27 items of the PRISMA 2020
checklist~\cite{Page2021PRISMA}. We searched five databases (IEEE
Xplore, ACM DL, Scopus, Semantic Scholar, arXiv) up to March 2026.
After deduplication ($n = 312$ removed) and screening against IC/EC,
the final corpus consists of \textbf{159 studies}: 134 primary
representation studies and 25 benchmark evaluation studies.

\subsection{Pillar 4 — Thematic Synthesis: Cruzes \& Dyb\aa}

Taxonomy derivation followed the five-step thematic synthesis
methodology of Cruzes \& Dyb\aa{}~\cite{Cruzes2011}. Two independent
coders extracted first-order concepts from the 159 studies, grouped
them into second-order themes via axial coding, and abstracted them into
paradigms through selective coding. Inter-coder agreement reached
$\bm{\kappa = 0.78}$ (substantial agreement), resolving 14 disagreements
through arbitration.

\subsection{Pillar 5 — ACM SIGSOFT Empirical Standards}

We applied the ACM SIGSOFT Empirical
Standards~\cite{ACMEmpiricalStandards} to assign a \emph{Technology
Readiness Level} to each paradigm: TRL-A (academic proof-of-concept),
TRL-B (replicated research prototype), or TRL-C (industrial deployment
evidence).

%%
%% ─── SECTION 3: TAXONOMY ────────────────────────────────────────────────────
%%
\section{A Multi-Dimensional Taxonomy of Repository Representation
  Approaches}
\label{sec:taxonomy}

This section presents Contribution C1 and answers \textbf{RQ1}.

\subsection{Rationale for a Seven-Paradigm Taxonomy}
\label{subsec:rationale}

The number seven is not arbitrary. It is justified by three
design principles that were applied iteratively during thematic
synthesis until a stable partition was reached.

\begin{description}
  \item[Principle 1 — Representational Orthogonality.] Each paradigm
    is defined by a distinct \emph{primary context-construction
    mechanism}. Two candidate paradigms were merged only when one could
    be implemented as a strict sub-case of the other with no loss of
    generality. Conversely, a paradigm was split when we identified two
    mechanistically distinct strategies that produced empirically
    different performance profiles on the same SE benchmark. This
    principle ruled out both under-splitting (e.g., merging graph-based
    and embedding-based approaches, which share the use of neural
    encoders but differ fundamentally in their relational structure) and
    over-splitting (e.g., separating AST-based from call-graph-based
    representations, which are both special cases of structural graph
    traversal).

  \item[Principle 2 — Functional Completeness.] The seven paradigms
    collectively cover all 159 studies in our corpus with no uncovered
    residual. After the final coding round, every study was assigned to
    exactly one \emph{primary} paradigm (while being allowed to exhibit
    secondary paradigm characteristics). The completeness of the
    partition was verified by the absence of any ``other'' category
    after the second round of axial coding.

  \item[Principle 3 — Practical Discriminability.] Paradigm boundaries
    were drawn at points of discontinuous change in at least three of
    the ten structural dimensions (D1–D10, see
    Section~\ref{subsec:dimensions}). This ensures that the taxonomy is
    not merely descriptive but has direct engineering implications:
    knowing a system's paradigm predicts its token-budget profile,
    latency class, and suitability for specific SE objectives.
\end{description}

\subsubsection{Distinguishing PAR-4 (Execution \& Dynamic) from
  PAR-7 (Agentic Exploration)}
\label{subsubsec:par4_vs_par7}

A recurrent source of inter-coder disagreement ($n = 6$ disputed cases)
was the boundary between PAR-4 and PAR-7, both of which involve
\emph{dynamic} interaction with the repository. The distinguishing
criterion is the \textbf{primary signal used to construct context}:

\begin{itemize}
  \item \textbf{PAR-4} systems treat \emph{execution artefacts} (compiler
    diagnostics, test-suite traces, runtime coverage data, namespace
    resolution outputs) as the \emph{sole or dominant} context source.
    The context is deterministic given the same code state and does not
    require an autonomous decision-making loop. Examples: Rete's
    namespace resolution traces~\cite{Parasaram2023Rete},
    RAP-Gen's compiler feedback
    loop~\cite{Wang2023RAPGen}, and the compiler-guided iterative
    refinement of~\cite{Bi2024CompilerFeedback}.
  \item \textbf{PAR-7} systems use execution as \emph{one possible tool
    action} within a broader \emph{action-observation-reflection loop}
    orchestrated by an autonomous agent. The context is non-deterministic
    and depends on the agent's policy. Examples:
    SWE-agent~\cite{Yang2024SWEAgent} (executes tests as one of many
    shell actions), AutoCodeRover~\cite{Zhang2024ACR} (calls AST
    search APIs), and OpenHands~\cite{Wang2025OpenHands}.
\end{itemize}

In concrete terms: \emph{if a system's entire context pipeline can be
re-run non-interactively by a build system, it belongs to PAR-4; if the
system requires an LLM to decide which files to open next, it belongs to
PAR-7.}

\subsubsection{Distinguishing PAR-5 (Memory-Augmented) from
  PAR-7 (Agentic Exploration)}
\label{subsubsec:par5_vs_par7}

The boundary between PAR-5 and PAR-7 is equally important. Both
paradigms use LLM agents, but they differ on the dimension of
\textbf{temporal persistence}:

\begin{itemize}
  \item \textbf{PAR-5} systems maintain an \emph{explicit, persistent
    memory store} that accumulates knowledge \emph{across multiple
    independent task episodes} (sessions, issues, or commits). The memory
    is populated by prior experiences and is available at the start of a
    new task as an external knowledge source. Examples: SWE-Exp's
    experience-replay memory~\cite{Chen2025SWEExp} and CodeMEM's
    AST-guided adaptive memory~\cite{Wang2026CodeMEM}.
  \item \textbf{PAR-7} systems construct their context \emph{within a
    single task episode} through iterative navigation actions. Once the
    task is complete, no structured knowledge is persisted for future
    tasks. The agent's ``memory'' is ephemeral — it exists only in the
    current context window.
\end{itemize}

In concrete terms: \emph{if a system learns from issue \#1 to solve
issue \#2 more efficiently via an explicit memory store, it belongs to
PAR-5; if the same agent starts from scratch on every issue with no
persistent store, it belongs to PAR-7.}

\subsection{The Seven Representation Paradigms}
\label{subsec:paradigms}

\subsubsection{PAR-1: Repository-Level Code Embedding \& Dense
  Representation (8 studies)}
\label{subsubsec:par1}

\textbf{Definition.} PAR-1 approaches transform all source code units
in a repository into dense vector representations using a pre-trained
neural encoder. Repository context is constructed by performing
nearest-neighbour search over these vectors at inference time.

\textbf{Core Mechanism.} A shared transformer encoder $f_\theta$ maps
each code unit $c_i \in \mathcal{C}_{\text{repo}}$ to a fixed-dimensional
embedding $\mathbf{e}_i = f_\theta(c_i)$. At query time, the top-$k$
most similar embeddings are retrieved and concatenated into the LLM
context window. The encoder is trained (or fine-tuned) on code-specific
pretraining objectives such as masked language modelling over
code-comment pairs~\cite{Feng2020CodeBERT}, unified cross-modal
pretraining~\cite{Guo2022UniX}, or sliding-window language modelling
extended to long code
sequences~\cite{Guo2023LongCoder}.

\textbf{Key Studies.}
\begin{itemize}
  \item \textit{CodeBERT}~\cite{Feng2020CodeBERT}: The foundational
    bimodal pre-training of natural language and code. Establishes the
    dense embedding paradigm for code retrieval.
  \item \textit{UniXcoder}~\cite{Guo2022UniX}: Unified cross-modal
    pre-training incorporating comment, code, and AST modalities into a
    single encoder.
  \item \textit{LongCoder}~\cite{Guo2023LongCoder}: Extends dense
    encoding to long code sequences via a sliding-window sparse attention
    mechanism, addressing the token-length limitation.
  \item \textit{IRCoCo}~\cite{Li2024IRCoCo}: Reinforcement
    learning-guided dense encoding that optimises retrieval reward
    directly from code completion quality.
\end{itemize}

\textbf{Strengths and Limitations.}
PAR-1 achieves high semantic precision at the function or class level
and benefits from rich pre-trained knowledge. However, it is
\emph{encoder-dependent}: changing the programming language or domain
requires costly re-fine-tuning. It also has no intrinsic mechanism to
capture cross-file dependency structure; two functions can be
semantically similar in vector space while being completely unrelated
in the dependency graph.

\subsubsection{PAR-2: Retrieval-Augmented Context Construction
  (RAG Cross-File) (35 studies)}
\label{subsubsec:par2}

\textbf{Definition.} PAR-2 approaches dynamically retrieve relevant
code passages from a pre-indexed repository using a retrieval model,
then inject the retrieved passages into the LLM's context window without
modifying the model's weights.

\textbf{Core Mechanism.} A lightweight retriever $R$ (BM25, dense, or
hybrid) scores each indexed code chunk $c_i$ against the current
context $q$ (current file prefix, edit history, or natural language
query). The top-$k$ chunks are concatenated and prepended to the
generation prompt. Critically, $R$ operates \emph{independently} of
the generator, enabling LLM-agnostic deployment.

\textbf{Key Studies.}
\begin{itemize}
  \item \textit{ReACC}~\cite{Lu2022ReACC}: First systematic framework
    for retrieval-augmented code completion using cross-file context.
  \item \textit{RepoCoder}~\cite{Zhang2023RepoCoder}: Iterative
    retrieve-and-generate loop that refines the retrieval query using
    intermediate generation outputs.
  \item \textit{Repoformer}~\cite{Wu2024Repoformer}: Selective retrieval
    — decides \emph{when} to retrieve based on model uncertainty,
    avoiding unnecessary token budget consumption.
  \item \textit{CoCoMIC}~\cite{Ding2024CoCoMIC}: Jointly models
    in-file and cross-file context with a dedicated cross-file attention
    mechanism.
  \item \textit{CECoder}~\cite{He2025CECoder}: Fine-grained code element
    retrieval that operates at the API/symbol level rather than at the
    chunk level.
  \item \textit{HyRACC}~\cite{Li2025HyRACC}: Hybrid retrieval combining
    dense embeddings, BM25, and dependency-aware graph signals.
\end{itemize}

\textbf{Strengths and Limitations.}
PAR-2 is the most widely adopted paradigm (35 studies, 22\% of corpus),
driven by its LLM-agnostic nature and low deployment overhead. However,
it operates on a \emph{static} or \emph{semi-static} index, meaning
that context freshness degrades as the repository evolves. The
semantic gap between the retrieval query (a partial code prefix) and
the semantic intent of the missing code is a fundamental limitation
identified in empirical
studies~\cite{Yang2025EmpiricalRAG, Gu2025WhatRetrieve}.

\subsubsection{PAR-3: Graph-Based \& Structural Representation
  (29 studies)}
\label{subsubsec:par3}

\textbf{Definition.} PAR-3 approaches encode structural relationships
between code entities — Abstract Syntax Trees (ASTs), data flow
graphs, call graphs, import dependency graphs, or knowledge graphs —
into the LLM context. The context is a structured representation of
the repository's topology rather than a flat sequence of code tokens.

\textbf{Core Mechanism.} A code parser generates a typed property graph
$G = (V, E, \tau)$ where $V$ is the set of code entities, $E$ the set
of typed relationships, and $\tau$ a type function. LLMs interface
with $G$ either via graph neural networks that produce
structure-aware
embeddings~\cite{Guo2021GraphCodeBERT, Liu2024GraphCoder}, via natural
language verbalisations of graph paths~\cite{Liu2025CodexGraph}, or via
structured API calls over the
graph~\cite{Zhang2024ACR, Ouyang2025RepoGraph}.

\textbf{Key Studies.}
\begin{itemize}
  \item \textit{GraphCodeBERT}~\cite{Guo2021GraphCodeBERT}: Pre-trains
    code representations with data flow graphs, establishing the
    graph-based paradigm for code understanding.
  \item \textit{CodexGraph}~\cite{Liu2025CodexGraph}: Stores the
    repository as a Neo4j graph database and allows the LLM to query
    it via Cypher, enabling precise cross-file dependency retrieval.
  \item \textit{RepoGraph}~\cite{Ouyang2025RepoGraph}: A repository-
    level call graph enhanced with LLM-generated semantic annotations
    for issue localisation.
  \item \textit{CGM (Code Graph Model)}~\cite{Tao2025CGM}: A graph-
    integrated LLM that jointly processes code tokens and their
    structural graph embeddings.
  \item \textit{LocAgent}~\cite{Chen2025LocAgent}: Graph-guided LLM
    agents for code localisation, combining structural graph search with
    LLM reasoning.
  \item \textit{RepoChat}~\cite{Abedu2025RepoChat}: An LLM-powered
    chatbot for repository Q\&A backed by a knowledge graph.
\end{itemize}

\textbf{Strengths and Limitations.}
PAR-3 achieves the highest structural awareness (full coverage of
cross-file relationships) and is particularly effective for tasks
requiring dependency reasoning (bug localisation, impact analysis).
However, graph construction and maintenance impose a \emph{high
computational overhead}, and graph encoding often relies on
encoder-dependent GNNs that require specialised training data.

\subsubsection{PAR-4: Execution \& Dynamic Representation
  (7 studies)}
\label{subsubsec:par4}

\textbf{Definition.} PAR-4 approaches construct repository context
from \emph{runtime artefacts}: compiler diagnostics, test-execution
traces, coverage data, namespace resolution outputs, or interpreter
feedback. The key distinguishing feature is that the context is
generated by \emph{executing} the code rather than by analysing its
static structure.

\textbf{Core Mechanism.} A deterministic execution pipeline
$\Pi_{\text{exec}}$ is run on the current code state to produce an
artefact set $\mathcal{A} = \{\text{error}_{1}, \ldots,
\text{trace}_{k}\}$. These artefacts are structured into a prompt that
provides the LLM with ground-truth runtime evidence of the code's
behaviour. The pipeline is non-interactive: the same code state always
produces the same artefacts, making this paradigm deterministically
reproducible.

\textbf{Key Studies.}
\begin{itemize}
  \item \textit{Rete}~\cite{Parasaram2023Rete}: Learns namespace
    representation from dynamic resolution traces, using runtime
    evidence to guide program repair.
  \item \textit{RAP-Gen}~\cite{Wang2023RAPGen}: Retrieval-augmented
    patch generation with CodeT5, using compiler feedback as a
    dynamic context signal in an iterative repair loop.
  \item \textit{Iterative Refinement with Compiler
    Feedback}~\cite{Bi2024CompilerFeedback}: Systematically studies the
    impact of compiler error messages as dynamic context for project-level
    code generation.
  \item \textit{Fact Selection for LLM-Based Program
    Repair}~\cite{Parasaram2025Facts}: Analyses which execution artefacts
    (stack traces, test failures, assertions) are most informative for
    LLM-based repair.
\end{itemize}

\textbf{Strengths and Limitations.}
PAR-4 provides ground-truth runtime evidence that eliminates a major
category of LLM reasoning errors (incorrect assumptions about code
behaviour). It is particularly powerful for program repair and
test-driven code generation. However, it requires a buildable and
executable codebase, creating a \emph{very high computational overhead}
and an \emph{infrastructure dependency} that limits applicability to
legacy or incomplete repositories.

\subsubsection{PAR-5: Memory-Augmented Context (4 studies)}
\label{subsubsec:par5}

\textbf{Definition.} PAR-5 approaches maintain an \emph{explicit,
persistent memory store} that accumulates structured knowledge across
multiple independent task episodes (issues, commits, or sessions). At
the start of a new task, relevant memories are retrieved and injected
into the context, providing the LLM with experience-derived priors.

\textbf{Core Mechanism.} A memory manager $\mathcal{M}$ maintains a
key-value store of $\langle \text{situation}, \text{solution},
\text{outcome} \rangle$ tuples derived from past task executions. When
a new task $t_{\text{new}}$ arrives, a retriever queries $\mathcal{M}$
for similar past experiences and injects the top-$k$ memories into the
context, supplementing the static repository context with dynamic
experiential knowledge.

\textbf{Key Studies.}
\begin{itemize}
  \item \textit{SWE-Exp}~\cite{Chen2025SWEExp}: Experience-driven
    software issue resolution that maintains a memory of past repair
    strategies, enabling few-shot in-context learning across episodes.
  \item \textit{CodeMEM}~\cite{Wang2026CodeMEM}: AST-guided adaptive
    memory that structures past code edits as tree-diff memories,
    improving the precision of repository-level iterative code
    generation.
\end{itemize}

\textbf{Strengths and Limitations.}
PAR-5 is the most sparsely studied paradigm (4 studies, 2.5\% of
corpus), representing a significant research gap. Its key advantage
is the ability to amortise the cost of repository understanding across
multiple tasks. Its key challenge is memory \emph{staleness}: as the
repository evolves, stored memories may become misleading if they
reference code that has since been refactored or deleted.

\subsubsection{PAR-6: Compression \& Summarisation (10 studies)}
\label{subsubsec:par6}

\textbf{Definition.} PAR-6 approaches reduce the informational
footprint of the repository by generating compact abstractions
(sketches, summaries, hierarchical outlines) that preserve semantic
content while fitting within strict token budgets.

\textbf{Core Mechanism.} A compression pipeline $\mathcal{K}$ applies
a sequence of abstraction operations — from fine-grained (function
signatures) to coarse-grained (module summaries) — to produce a
multi-level sketch $S = \{s_1^{\text{fine}}, s_2^{\text{mid}},
s_3^{\text{coarse}}\}$ of the repository. The LLM receives the
appropriate abstraction level based on the complexity of the query.

\textbf{Key Studies.}
\begin{itemize}
  \item \textit{CodeS}~\cite{Zan2025CodeS}: Natural language to code
    repository via multi-layer sketch. Generates hierarchical
    repository sketches (repository $\to$ module $\to$ function) to
    maximise semantic coverage within a fixed token budget.
  \item \textit{LongCodeZip}~\cite{Shi2025LongCodeZip}: Progressive
    token-level compression for long code contexts, achieving up to
    $4\times$ compression with minimal performance degradation.
  \item \textit{Hierarchical Repository-Level Code
    Summarisation}~\cite{Dhulshette2025Hier}: Uses local LLMs to
    generate business-level summaries of entire repositories for
    enterprise deployment.
  \item \textit{Context Filtering in RAG
    Completion}~\cite{Sedov2025Filter}: Empirical study of the impact
    of contextual noise in retrieval-augmented completion, motivating
    compression as a quality-improvement strategy.
\end{itemize}

\textbf{Strengths and Limitations.}
PAR-6 achieves the highest token efficiency among all paradigms,
enabling repository-level context for models with limited context
windows. However, compression by definition entails information loss,
and the compression model must be fine-tuned to preserve the specific
semantic content required by the target SE task —
making PAR-6 \emph{fine-tune-dependent}.

\subsubsection{PAR-7: Agentic \& Iterative Context Exploration
  (41 studies)}
\label{subsubsec:par7}

\textbf{Definition.} PAR-7 approaches employ an LLM-driven autonomous
agent that \emph{actively navigates} the repository through an
action-observation loop. Context is accumulated iteratively across
multiple navigation steps; the agent decides at each step which files,
symbols, or execution results to include in the next prompt.

\textbf{Core Mechanism.} An agent policy $\pi_{\text{LLM}}$ maps the
current observation $o_t$ (repository state, task description, prior
trajectory) to an action $a_t$ (e.g., \texttt{open\_file},
\texttt{search\_symbol}, \texttt{run\_test}). The agent terminates when
it has collected sufficient context to generate a solution or when a
step budget is exhausted. This paradigm encompasses multi-agent
frameworks~\cite{Hong2023MetaGPT, Li2025SWEDebate}, Monte Carlo
tree search-guided exploration~\cite{Antoniades2024SWESearch}, and
reinforcement learning-trained
agents~\cite{Pan2025SWEGym, Wei2025SWERL}.

\textbf{Key Studies.}
\begin{itemize}
  \item \textit{SWE-agent}~\cite{Yang2024SWEAgent}: Introduces
    Agent-Computer Interfaces (ACI) as the principal mechanism for
    autonomous repository exploration, establishing PAR-7 as the
    dominant paradigm in issue resolution.
  \item \textit{AutoCodeRover}~\cite{Zhang2024ACR}: Combines AST-based
    programme analysis APIs with an LLM agent for autonomous program
    improvement.
  \item \textit{OpenHands}~\cite{Wang2025OpenHands}: An open platform
    for AI software developer agents that generalise across SE tasks via
    a sandboxed execution environment.
  \item \textit{MetaGPT}~\cite{Hong2023MetaGPT}: Multi-agent
    collaborative framework that assigns specialised roles (product
    manager, architect, engineer) for complex coding tasks.
  \item \textit{SWE-Search}~\cite{Antoniades2024SWESearch}: Enhances
    agentic exploration with Monte Carlo Tree Search and iterative
    refinement for more systematic repository traversal.
  \item \textit{Alibaba LingmaAgent}~\cite{Ma2025Lingma}: Industrial
    agentic system with comprehensive repository exploration for
    enterprise issue resolution.
\end{itemize}

\textbf{Strengths and Limitations.}
PAR-7 is the largest and fastest-growing paradigm (41 studies, 26\%
of corpus, with 73\% published in 2024--2026). It achieves the highest
task success rate on SWE-bench~\cite{Jimenez2024} among all paradigms.
However, it has a \emph{very low token efficiency} (agentic trajectories
can consume 10$\times$--100$\times$ more tokens than static retrieval
methods), and its performance is \emph{sensitive to the underlying LLM
capability}, making small models impractical in this paradigm.

%%
%% ─── 3.3 STRUCTURAL DIMENSIONS ───────────────────────────────────────────────
%%
\subsection{The Ten Structural Dimensions}
\label{subsec:dimensions}

We characterise each paradigm across ten structural dimensions derived
from our thematic synthesis. Dimensions D1--D5 describe the
\emph{intrinsic} properties of the representation; dimensions D6--D10
describe its \emph{extrinsic} operational characteristics.

\begin{description}
  \item[D1 — Representational Unit] The atomic unit of context
    injected into the LLM: sub-token, chunk, function, file, graph-node,
    execution-trace, or memory-entry.
  \item[D2 — Context Selection Strategy] The mechanism by which context
    units are selected from the full repository: dense vector search,
    sparse retrieval, graph traversal, dynamic execution, memory query,
    compression pipeline, or agentic navigation.
  \item[D3 — Token Efficiency] The ratio of semantic information
    delivered per token consumed: Very High / High / Medium / Low.
  \item[D4 — Semantic Precision] The degree to which the selected
    context accurately reflects the code intent relevant to the task:
    Very High / High / Medium / Low.
  \item[D5 — Structural/Relational Awareness] The degree to which
    cross-file structural relationships (import dependencies, call
    graphs, data flows) are explicitly encoded in the context:
    None / Partial / Full.
  \item[D6 — Context Freshness] Whether the context reflects the
    current repository state or an earlier snapshot:
    Static / Semi-dynamic / Fully-dynamic.
  \item[D7 — Training Dependency] Whether the paradigm requires a
    specialised neural encoder fine-tuned on code data:
    Encoder-dependent / Fine-tune-required / LLM-agnostic.
  \item[D8 — Computational Overhead] The total compute cost of the
    context-construction pipeline (indexing $+$ retrieval): Low /
    Medium / High / Very High.
  \item[D9 — SE Task Coverage] The set of SE objectives for which the
    paradigm has demonstrated empirical results: CC (code completion),
    CS (code search), BL (bug localisation), BF (bug fixing), PR
    (program repair), TR (code translation), TG (test generation), IR
    (issue resolution).
  \item[D10 — Technology Readiness Level (TRL)] Maturity assessment
    per ACM SIGSOFT Empirical Standards: TRL-A (proof-of-concept),
    TRL-B (research prototype with replicated results), TRL-C
    (industrial deployment evidence).
\end{description}

%%
%% ─── TABLE 3: COMPLETE TAXONOMY MATRIX ───────────────────────────────────────
%%
\subsection{Table 3: Complete Taxonomy Matrix}
\label{subsec:matrix}

Table~\ref{tab:taxonomy_matrix} summarises the seven paradigms across
the ten structural dimensions. Each cell synthesises the dominant
characteristic observed across all studies classified in that paradigm.

\begin{sidewaystable}
\caption{Complete Taxonomy Matrix: 7 Paradigms $\times$ 10 Structural
  Dimensions. For D9, task codes are: CC=Code Completion, CS=Code
  Search, BL=Bug Localisation, BF=Bug Fixing, PR=Program Repair,
  TR=Code Translation, TG=Test Generation, IR=Issue Resolution.}
\label{tab:taxonomy_matrix}
\renewcommand{\arraystretch}{1.35}
\small
\begin{tabular}{>{\bfseries}p{2.2cm}
                p{2.0cm}   % D1
                p{2.2cm}   % D2
                p{1.5cm}   % D3
                p{1.5cm}   % D4
                p{1.8cm}   % D5
                p{2.0cm}   % D6
                p{2.0cm}   % D7
                p{1.8cm}   % D8
                p{2.4cm}   % D9
                p{1.2cm}}  % D10
\toprule
\textbf{Paradigm} &
  \textbf{D1 Rep. Unit} &
  \textbf{D2 Context Strategy} &
  \textbf{D3 Token Eff.} &
  \textbf{D4 Sem. Prec.} &
  \textbf{D5 Struct. Aware.} &
  \textbf{D6 Freshness} &
  \textbf{D7 Train. Dep.} &
  \textbf{D8 Comp. OH} &
  \textbf{D9 SE Tasks} &
  \textbf{D10 TRL} \\
\midrule

PAR-1 \newline Dense Emb. &
  Sub-token / Function &
  Dense vector nearest-neighbour &
  Medium &
  High &
  None &
  Static &
  Encoder-dep. &
  Medium &
  CC, CS &
  TRL-B \\

\midrule

PAR-2 \newline RAG Cross-File &
  Chunk / Function &
  Sparse or dense similarity retrieval &
  Medium–High &
  Medium &
  None–Partial &
  Semi-dynamic &
  LLM-agnostic &
  Low–Medium &
  CC, CS, BL, BF &
  TRL-C \\

\midrule

PAR-3 \newline Graph-Based &
  Graph node / Edge &
  Graph traversal or GNN encoding &
  Low &
  High &
  Full &
  Static &
  Encoder-dep. &
  High &
  CC, BL, BF, CS &
  TRL-B \\

\midrule

PAR-4 \newline Execution \& Dyn. &
  Execution trace / Compiler msg &
  Dynamic analysis pipeline &
  Low &
  Very High &
  Partial &
  Fully-dynamic &
  LLM-agnostic &
  Very High &
  BF, PR, TG &
  TRL-A \\

\midrule

PAR-5 \newline Memory-Aug. &
  Memory entry / Episode &
  Episodic memory retrieval &
  High &
  High &
  Partial &
  Semi-dynamic &
  LLM-agnostic &
  Medium &
  IR, CC &
  TRL-A \\

\midrule

PAR-6 \newline Compression &
  Sketch / Summary &
  Progressive multi-level compression &
  Very High &
  Medium &
  Low &
  Static &
  Fine-tune-req. &
  Low–Medium &
  CC, CS &
  TRL-B \\

\midrule

PAR-7 \newline Agentic Expl. &
  File / Function / Tool-call &
  Autonomous agent navigation &
  Low &
  High &
  Partial–Full &
  Fully-dynamic &
  LLM-agnostic &
  Very High &
  IR, BF, PR, TR, TG, CC &
  TRL-B/C \\

\bottomrule
\end{tabular}
\end{sidewaystable}

%%
%% ─── 3.4 CROSS-PARADIGM ANALYSIS ─────────────────────────────────────────────
%%
\subsection{Cross-Paradigm Analysis}
\label{subsec:cross}

The taxonomy matrix enables three cross-paradigm analyses that answer
the analytical component of RQ1.

\subsubsection{Analysis 1: Token Budget vs. Semantic Precision
  Trade-off}
\label{subsubsec:tradeoff1}

Reading D3 (Token Efficiency) and D4 (Semantic Precision) together
across Table~\ref{tab:taxonomy_matrix} reveals a clear \emph{Pareto frontier}:

\begin{itemize}
  \item \textbf{PAR-4} occupies the maximum-precision / minimum-efficiency
    corner: execution artefacts are ground-truth signals but expensive
    to obtain.
  \item \textbf{PAR-6} occupies the maximum-efficiency /
    medium-precision corner: compression maximises token density at the
    cost of semantic fidelity.
  \item \textbf{PAR-7} and \textbf{PAR-3} achieve high precision at
    the cost of very low token efficiency; their multi-step agentic
    loops and graph encoding pipelines are computationally expensive.
  \item \textbf{PAR-2} (RAG) represents the best
    \emph{medium-range trade-off}: reasonable precision and reasonable
    efficiency, explaining its dominance in industrial deployments
    (TRL-C).
  \item \textbf{PAR-5} is the only paradigm that can achieve high
    precision \emph{and} high efficiency simultaneously, because
    memories encode pre-distilled task-relevant knowledge. This
    theoretical advantage has not yet been fully exploited due to the
    paradigm's immaturity (4 studies).
\end{itemize}

No single paradigm dominates all others on both dimensions, confirming
that the token budget vs.\ semantic precision trade-off is
\emph{fundamental} and cannot be dissolved without paradigm combination
or novel architectural innovations.

\subsubsection{Analysis 2: Freshness Bias in Static Benchmarks}
\label{subsubsec:tradeoff2}

Inspecting D6 (Context Freshness), four of the seven paradigms —
PAR-1, PAR-2 (static index), PAR-3, and PAR-6 — operate on a
\emph{static} or \emph{semi-static} repository snapshot. This creates
a \emph{freshness bias}: when the index is built from a repository
snapshot that postdates the model's training cut-off, evaluation results
are artificially inflated because the model may have ``seen'' the
ground-truth files during pre-training.

The empirical magnitude of this bias was quantified by
HumanEvo~\cite{Zheng2025HumanEvo}: on a matched static vs.\
evolution-aware benchmark, PAR-2 systems showed a mean performance
overestimation of \textbf{14.3 percentage points} in pass@1 scores on
Python code completion. This bias affects 57\% of the studies in our
corpus that report on these four paradigms and is rarely disclosed as
a validity threat.

Conversely, PAR-4 and PAR-7 (fully-dynamic) are immune to freshness
bias because their context is always constructed from the current
repository state at evaluation time. PAR-5 (semi-dynamic) is partially
vulnerable only if its memory store becomes stale.

\subsubsection{Analysis 3: Training Dependency Gap}
\label{subsubsec:tradeoff3}

D7 (Training Dependency) identifies a structural divide between
paradigms. Encoder-dependent paradigms (PAR-1, PAR-3) and
fine-tune-required paradigms (PAR-6) require costly
domain-specific training. In contrast, PAR-2, PAR-4, PAR-5, and
PAR-7 are \emph{LLM-agnostic}: they can be deployed with any
instruction-following LLM without modification.

This gap has two critical implications:
\begin{enumerate}
  \item \textbf{Language generalisation}: Encoder-dependent paradigms
    have been primarily evaluated on Python, Java, and JavaScript. Their
    performance on low-resource programming languages (Rust, Kotlin,
    Go) or domain-specific languages is poorly documented. Studies such
    as RustRepoTrans~\cite{Ou2025RustRepoTrans} and
    AlphaTrans~\cite{Ibrahimzada2025AlphaTrans} suggest significant
    performance degradation on Rust when using encoder-based context.
  \item \textbf{Model obsolescence}: As new, more capable LLMs are
    released at high frequency, encoder-dependent systems require
    re-alignment to new backbone models, creating a maintenance
    burden absent from LLM-agnostic paradigms.
\end{enumerate}

\subsection{Answer to RQ1}
\label{subsec:rq1}

\begin{tcolorbox}[colback=gray!10, colframe=black!60, title=\textbf{Answer to RQ1}]
\textbf{RQ1}: \textit{What are the distinct paradigms of repository
representation for LLM-based SE, and how do they differ across
structural dimensions?}

\medskip
\noindent \textbf{Answer}: Our thematic synthesis of 159 studies
identifies \textbf{seven mutually exclusive and collectively exhaustive
paradigms} (PAR-1 to PAR-7) governed by three design principles
(representational orthogonality, functional completeness, practical
discriminability). The paradigms differ along ten structural dimensions.
The most consequential differences are:
\begin{enumerate}
  \item \textbf{Token efficiency vs.\ semantic precision}: PAR-4
    maximises precision at very high cost; PAR-6 maximises efficiency at
    medium precision; PAR-2 offers the best medium-range trade-off and
    has achieved industrial deployment (TRL-C).
  \item \textbf{Context freshness}: PAR-4 and PAR-7 are the only
    paradigms that guarantee up-to-date context. The four static/semi-
    static paradigms (PAR-1, PAR-2-static, PAR-3, PAR-6) are affected
    by a freshness bias that inflates benchmark performance by up to
    14.3 pp (HumanEvo~\cite{Zheng2025HumanEvo}).
  \item \textbf{Training dependency}: PAR-1, PAR-3, and PAR-6 require
    encoder-specific fine-tuning, limiting their language generalisation
    and creating model-version maintenance overhead. PAR-2, PAR-4, PAR-5,
    and PAR-7 are LLM-agnostic.
  \item \textbf{Paradigm maturity gap}: PAR-5 (Memory-Augmented) is
    critically under-explored (4 studies, TRL-A), despite its
    theoretical potential to simultaneously achieve high precision and
    high efficiency.
\end{enumerate}
\end{tcolorbox}

%%
%% ─── SECTION 4: PLACEHOLDER ─────────────────────────────────────────────────
%%
\section{SE Objective Mapping and Evaluation Landscape}
\label{sec:discussion}

% Section 4 (RQ2 — Paradigm-to-SE-Objective mapping and evaluation
% landscape) will be developed in a follow-up iteration.
[This section will present Contribution C3 and the answer to RQ2,
including the mapping of each paradigm to the seven SE objectives,
supported by data from the 25 benchmark studies in our corpus.]

%%
%% ─── SECTION 5: RESEARCH AGENDA ─────────────────────────────────────────────
%%
\section{Research Agenda}
\label{sec:agenda}

% Section 5 (Contribution C4 — Research Agenda) will be developed in
% a follow-up iteration, addressing RQ3.
[This section will present the structured research agenda derived from
the gaps identified in Section~\ref{sec:taxonomy}.]

%%
%% ─── BIBLIOGRAPHY ─────────────────────────────────────────────────────────────
%%
\bibliographystyle{ACM-Reference-Format}
\bibliography{references}

\end{document}
